\section{First proof attacks: exact lemmas to establish or break}

The following statements turn the review's broad objections into falsifiable mathematical targets. A proof must derive the maps from the analytic charts; matching ranks is insufficient.

\subsection{Cusp quotient and collapse}

The source defines an infinite toric threefold $Y$ from the height-one cones over the periodic $A_2$ triangulation and an action of the form
\[
\Psi_{\bar\lambda}(x,t)
=\bigl(c_{\bar\lambda}(t)t^{B_0\bar\lambda}x,t\bigr),
\qquad \bar\lambda\in\bar\Lambda.
\]
Theorem~4.5's analytic part has now been independently retraced: this action
is free and properly discontinuous, its quotient is a smooth Hausdorff
threefold proper over the disc, and its central fibre has the asserted local
normal-crossings equations.  C1 below is narrower: it asks for the finite
equivariant topological model and collapse that are not supplied by those
analytic quotient facts.

\begin{obligation}[C1: finite equivariant cusp model]\label[obligation]{obl:C1}
Construct a locally finite $\bar\Lambda$-equivariant CW model of $(Y_{\le\eta},Y_0)$ whose quotient is finite, and prove that the collapse is equivariant in the diagram
\[
\begin{CD}
Y_{\le\eta} @>{r_Y}>> Y_0\\
@VV{/\bar\Lambda}V @VV{/\bar\Lambda}V\\
N'_0 @>{r}>> W.
\end{CD}
\]
The proof may cite the verified analytic action theorem, but it must still
establish continuity of the proposed collapse across the one-, two-, and
three-branch strata, compatibility with the fibre-torus action, finiteness of
the quotient CW model, and a strong deformation retraction.
\end{obligation}

\begin{obligation}[C2: integral collapse matrix]\label[obligation]{obl:C2}
With fixed cellular orientations, prove for $0\le q\le4$ the exact sequence
\[
0\longrightarrow
\operatorname{im}(\wedge^qM_0-I)
\longrightarrow\wedge^q\Lambda
\xrightarrow{c_*}H_q(W;\mathbb Z)
\longrightarrow0.
\]
In particular, derive $c_*\colon\Lambda\to H_1(W)=\bar\Lambda$ as the quotient map, prove
\[
c_*(\widehat\gamma\wedge\widehat u)=[\mathbb T_B],
\qquad
c_*(\widehat w\wedge\widehat\delta)=0,
\]
and record the degree-three orientations rather than inferring them from ranks.
\end{obligation}

\begin{obligation}[C3: collapse equals specialization]\label[obligation]{obl:C3}
Prove functorially and integrally that
\[
c^*=i_F^*r^*=\operatorname{sp}_{\mathrm{nc}}^q,
\qquad
\operatorname{im}c^*=(\wedge^qV)^{T_0},
\]
including saturation, torsion, bases, and the degree-three sign. The comparison diagram must state the precise nearby-cycle edge map; the equality may not be defined by choosing the manuscript's answer as a normalization.
\end{obligation}

\subsection{Based inclusions}

Choose one basepoint on the zero section and fixed whiskers through every collar.

\begin{obligation}[V1: primitive meridians and drift]\label[obligation]{obl:V1}
Derive from the coordinate gluings, with no unrecorded conjugations,
\[
\widehat g_j^{m_j}=t_{v_j},
\qquad
\widetilde a_0=t_\mu\widehat\rho_2^{-1}\widehat\rho_1^{-1}.
\]
Prove that the clockwise meridian lifts to $\widehat g_j$ rather than its inverse or a nonprimitive power, that the two finite-fibre signs agree with the chosen orientations, that the zero-section whiskers add no hidden translation, and that the tautological cusp disc gives $\mu=0$.
\end{obligation}

\subsection{Mayer--Vietoris and Leray}

\begin{obligation}[MV1: geometric boundary cylinders]\label[obligation]{obl:MV1}
Using the independently constructed collars and the maps from C2--C3, prove
for every relevant degree $q$ that
\[
 \alpha_q=
 \bigl[(i_{U*},-i_{N*}):H_q(U\cap N;\mathbb Z)
 \longrightarrow H_q(U;\mathbb Z)\oplus H_q(N;\mathbb Z)\bigr]_{\mathcal B_q}
\]
in the manuscript's named primitive oriented bases.  This must include a
chain-level check of every zero column, sign, torsion class, and saturation
claim; equality of ranks or Smith forms is not sufficient.  For
$K_k=\ker(H_k(M)\to H_k(W))$, derive geometrically
\[
K_1=\mathbb Z[\widetilde a_0],
\qquad
K_2=\langle\xi_{\widehat w},\xi_{\widehat\delta}\rangle,
\]
and fix one orientation convention for
\[
\partial(j_*\xi_\nu)=(I-A_1^{-1})\nu=(I-A_2)\nu.
\]
The same calculation must recover the multiple-fibre class
\([F]\longmapsto(3[S_1],-4[S_2])\) and the index-two degree-three image at $p_2$.
\end{obligation}

\begin{obligation}[L1: transgression comparison]\label[obligation]{obl:L1}
Construct the clutching cocycle that compares the Leray transgression with the Mayer--Vietoris obstruction. For $12\gamma$, derive the discrepancy triple
\[
\bigl(4\ell_1-\alpha_1,\ 3\ell_2-\alpha_2,\
-(12\ell_0-\alpha_1-\alpha_2)\bigr),
\]
whose sum is $-p$, and thereby determine---rather than assume---the sign in
\[
d_2^{0,1}(12\gamma)=\pm p\,\omega.
\]
Then prove that the cup-product and coinvariant identifications are natural in the same orientations before using multiplicativity in higher degrees.
\end{obligation}

Until C1--C3, V1, MV1, and L1 are established, the conclusion $X\cong S^6$ remains conditional. The integer algebra downstream of these statements is already fixed by the previous section.
